\begin{statementbox}
	\textbf{\textcolor{CStatement}{条件}}
	\begin{description}[style=nextline, leftmargin=2.8em, labelsep=0.5em, itemsep=0.35em]
		\item[\textbf{\textcolor{CStatement}{条件 1（线性约束）：}}]
			约束条件为线性不等式组，且可行域 $D$ 是有界凸多边形（或虽无界但最优解存在）。
		\item[\textbf{\textcolor{CStatement}{条件 2（线性目标）：}}]
			目标函数为 $z=ax+by$，其中 $a,b$ 不全为零。
	\end{description}
	\textbf{\textcolor{CStatement}{结论}}
	\begin{description}[style=nextline, leftmargin=2.8em, labelsep=0.5em, itemsep=0.45em]
		\item[\textbf{\textcolor{CStatement}{结论一（顶点最值）：}}]
			\textbf{\textcolor{CStatement}{直观描述：}}
			最值只会在可行域的顶点处出现，不会在边上内部点。
			\[
				\max_{D} z = \max_{\mathcal{V}} z(P),
				\qquad
				\min_{D} z = \min_{\mathcal{V}} z(P),
			\]
			其中 $\mathcal{V}$ 是 $D$ 的所有顶点构成的集合。
		\item[\textbf{\textcolor{CStatement}{推广结论（截距判别）：}}]
			\textbf{\textcolor{CStatement}{直观描述：}}
			通过等值线的截距正负判断最值方向。 当 $b\neq 0$ 时，目标函数化为 $y = -\dfrac{a}{b}x + \dfrac{z}{b}$。 若 $b>0$，则 $z$ 的最大值对应 $\dfrac{z}{b}$ 最大，最小值对应 $\dfrac{z}{b}$ 最小；若 $b<0$，则 $z$ 的最大值对应 $\dfrac{z}{b}$ 最小，最小值对应 $\dfrac{z}{b}$ 最大。
	\end{description}

\end{statementbox}
